\chapter*{Conclusiones}

El desarrollo de la aplicación web integrada al Sistema Integral de Producción permitió atender las necesidades operativas identificadas dentro de las áreas de producción y ensamblaje de TenarisTamsa, contribuyendo a la modernización y digitalización de procesos relacionados con la gestión de información productiva. A través de la implementación de módulos especializados, se logró fortalecer funcionalidades clave del sistema, mejorando la captura, consulta y administración de datos dentro del entorno operativo.

Durante el desarrollo del proyecto se aplicaron conocimientos relacionados con arquitectura de software, desarrollo frontend y backend, diseño de interfaces, integración de servicios y gestión de bases de datos, utilizando tecnologías modernas como React.js, .NET y PostgreSQL. La implementación de una arquitectura basada en componentes para el frontend y una arquitectura multicapa para el backend permitió desarrollar una solución modular, escalable y mantenible, alineada con los estándares tecnológicos definidos por la organización.

Asimismo, el sistema desarrollado permitió reducir la dependencia de registros manuales, facilitando la disponibilidad de información en tiempo real y mejorando la trazabilidad de los procesos productivos. La centralización de la información y la integración de distintos módulos dentro de una misma plataforma contribuyen a optimizar el seguimiento de órdenes de producción, el control operativo y la toma de decisiones dentro del área de trabajo.

De igual manera, el proyecto permitió reforzar la importancia de las buenas prácticas de desarrollo, documentación y estandarización del código en entornos empresariales e industriales de gran escala. La experiencia obtenida durante el periodo de residencias profesionales favoreció el desarrollo de competencias técnicas y profesionales relacionadas con el trabajo colaborativo, el análisis de requerimientos, la resolución de problemas y la adaptación a procesos corporativos de desarrollo de software.

Finalmente, se concluye que la implementación de soluciones tecnológicas orientadas a la digitalización y automatización de procesos representa un factor clave para mejorar la eficiencia operativa, la confiabilidad de la información y la capacidad de crecimiento de las organizaciones industriales; así como la reducción de la pérdida de información en un 35\% y la disminución del consumo de papel en un 10\%. El sistema desarrollado establece una base sólida para futuras mejoras y ampliaciones dentro del Sistema Integral de Producción, contribuyendo al fortalecimiento tecnológico y operativo de la empresa.